% 自 notes/docs/1151/seminar/W09_1109_預習.md（87808f3）由 tools/md_to_tex.py 於 2026-11-09 轉出一次；此後本檔為正本。
\chapter{Alexander 與 Schauer〈Defending Judicial Supremacy: A Reply〉\pg{455 至 482 頁}}
\label{ch:AS}

\begin{longtable}[]{@{}
  >{\raggedright\arraybackslash}p{(\linewidth - 4\tabcolsep) * \real{0.1200}}
  >{\raggedright\arraybackslash}p{(\linewidth - 4\tabcolsep) * \real{0.7200}}
  >{\raggedright\arraybackslash}p{(\linewidth - 4\tabcolsep) * \real{0.1600}}@{}}
\caption{Alexander 與 Schauer 一文的節次與頁次。}\tabularnewline
\toprule\noalign{}
\begin{minipage}[b]{\linewidth}\raggedright
節次
\end{minipage} & \begin{minipage}[b]{\linewidth}\raggedright
小標
\end{minipage} & \begin{minipage}[b]{\linewidth}\raggedright
頁次
\end{minipage} \\
\midrule\noalign{}
\endfirsthead
\toprule\noalign{}
\begin{minipage}[b]{\linewidth}\raggedright
節次
\end{minipage} & \begin{minipage}[b]{\linewidth}\raggedright
小標
\end{minipage} & \begin{minipage}[b]{\linewidth}\raggedright
頁次
\end{minipage} \\
\midrule\noalign{}
\endhead
\bottomrule\noalign{}
\endlastfoot
導言 & （無小標） & 455--458 \\
I & The Nature of the Question（A. The Preconstitutional Constitution；B. The Nature of the Inquiry） & 459--467 \\
II & The Settlement Function and Its Optimal Location（A. Is Settlement Overrated?；B. Are Courts the Best Settlers?） & 467--478 \\
III & What Is the Supreme Court For? & 478--481 \\
結論 & Conclusion & 481--482 \\
\end{longtable}

\section{導言：先前論點、批評與本文目的\pg{455 至 458 頁}}
\label{sec:AS-1}

\subsection{（一）原文的主張\pg{455 頁，註 1--4}}

Alexander\ix{persons}{alexander@Alexander, Larry} 與 Schauer\ix{persons}{schauer@Schauer, Frederick} 以 1997 年發表於 \en{Harvard Law Review} 的 On Extrajudicial Constitutional Interpretation 為前文（註 1）。本文開頭重述前文立場：最高法院對憲法的解釋，應被所有其他官員，不論司法或非司法，視為與憲法本身具有相同的權威地位\pg{455 頁}。兩人引 Hughes\ix{persons}{hughes@Hughes, Charles Evans} 的話（註 2）「the Constitution is what the judges say it is」〔憲法就是法官所說的那樣〕，承認這句話作為說明對最高法院論辯意味著什麼的嘗試，可能是糟糕的法理學，因為它無從理解（註 3 引 Dworkin\ix{persons}{dworkin@Dworkin, Ronald（德沃金）}, \en{Law\textquotesingle s Empire} 與 Hart\ix{persons}{hart@Hart, H. L. A.（哈特）}, \en{The Concept of Law}）。他們仍主張，這是非司法官員對最高法院及其見解宜有的態度，其情形與下級法院法官對最高法院及其意見宜有的態度相同，只是爭議小得多（註 4 引 Caminker 與 Schauer 論先例的著作）。

\subsection{（二）批評的來源與內容\pg{456 至 457 頁，註 5--9}}

這一立場通常與最高法院的 Cooper v. Aaron\ix{cases}{02@美國!01@聯邦最高法院!1958Cooper v. Aaron（1958）@Cooper v. Aaron（1958）}（1958；註 5）及較近的 City of Boerne\ix{cases}{02@美國!01@聯邦最高法院!1997City of Boerne v. Flores（1997）@City of Boerne v. Flores（1997）} v. Flores（1997；註 6，引 521 U.S. at 536）相連，長期受到強烈批評\pg{456 頁}。批評一部分來自學者（註 7 列 Levinson\ix{persons}{levinson@Levinson, Sanford}、Gant、McConnell、Nagel\ix{persons}{nagel@Nagel, Robert F.}、Paulsen 與 Tushnet\ix{persons}{tushnet@Tushnet, Mark} 1992 年的 The Constitution Outside the Courts），更多來自行政與立法部門的官員，後者看不出自己對憲法的解釋為何應因司法部門的先行主張而被推翻，而憲法文本並未賦予司法部門解釋上的優先地位（註 8 引 Meese\ix{persons}{meese@Meese, Edwin, III} 1987 年的 The Law of the Constitution，並指出 Lincoln\ix{persons}{lincoln@Lincoln, Abraham（林肯）} 對 Dred Scott\ix{cases}{02@美國!01@聯邦最高法院!1857Dred Scott v. Sandford（1857）@Dred Scott v. Sandford（1857）} 案的態度基於實質相同的立場）。批評者堅持，司法至上\ix{terms}{01@\zhlidxcat{司法審查權限之歸屬}!司法至上}的主張缺乏歷史、憲法體制結構與規範性的政治理論或制度設計論證的支持\pg{456 至 457 頁，註 9 引 Tribe\ix{persons}{tribe@Tribe, Laurence H.} 對 Cooper 的狹義解讀}。

註 6 另指出 City of Boerne 與本文的主要問題「orthogonal」〔正交〕：該案討論國會依增修條文第 14 條\ix{statutes}{02@美國!01@聯邦憲法!114@增修條文第 14 條}第 5 項的權力在何種情形下得偏離最高法院先例，屬實體解釋問題，兩人說這與服從問題有所不同。兩人同時表示，Kennedy\ix{persons}{kennedy@Kennedy, Anthony} 大法官所稱國會「has not just the right but the duty to make its own informed judgment on the meaning and force of the Constitution」〔不僅有權，而且有義務對憲法的意義與效力作出自己知情的判斷〕，是該意見中令人困擾的部分，他們認為那只是對國會的客氣，相當於稱讚敗訴律師的論證。

\subsection{（三）定紛止爭論的要點\pg{457 至 458 頁，註 10--11}}

最高法院在 Cooper\ix{cases}{02@美國!01@聯邦最高法院!1958Cooper v. Aaron（1958）@Cooper v. Aaron（1958）} 與 City of Boerne\ix{cases}{02@美國!01@聯邦最高法院!1997City of Boerne v. Flores（1997）@City of Boerne v. Flores（1997）} 宣告自身的至上地位，在兩人看來與問題無關，因為這種宣告預設了其他部門與公眾是否應接受該宣告的結論\pg{457 頁，註 10 引 Lawson \& Moore}。不依賴這種循環論證，兩人的前文主張可歸納為以下環節：

\begin{enumerate}
\def\labelenumi{\arabic{enumi}.}
\tightlist
\item
  法律的核心道德功能之一，是就應為之事作出定紛止爭\ix{terms}{02@\zhlidxcat{定紛止爭論}!定紛止爭}。
\item
  憲法因修改困難，比制定法與普通法裁判試圖建立更持久的定紛。
\item
  定紛的意義本身有爭議時，司法解釋的核心道德功能就在於確定該（企圖中的）定紛之意義。
\item
  司法解釋可能有誤，並不使其權威受到質疑，因為所有法律上的定紛都主張權威，即使受其拘束者認為它在道德上與法律上有誤。註 11 說明，主張權威即主張提供 content-independent\ix{terms}{02@\zhlidxcat{定紛止爭論}!與內容無關之理由} reasons〔與內容無關之理由〕，其成立與否取決於長期而言受拘束者以定紛為準，與逐案依自己的實體判斷行事相比，所犯錯誤的多寡\pg{註 11 引 Hart\ix{persons}{hart@Hart, H. L. A.（哈特）}、Raz 與 Schauer\ix{persons}{schauer@Schauer, Frederick} \en{Playing By the Rules} 128--34 頁}。
\item
  對第一階\ix{terms}{02@\zhlidxcat{定紛止爭論}!一階定紛與二階定紛}定紛的意義所作的錯誤第二階定紛，其權威的問題比第一階定紛本身的權威問題較小。因此期待人們即使在道德上不同意仍遵守憲法，便可基於同樣理由，期待他們在認為最高法院的解釋有誤時仍遵守之\pg{457 至 458 頁}。
\end{enumerate}

\subsection{（四）本文的目的\pg{458 頁，註 12--13}}

兩人說前文過去三年成為許多批評者的標的。註 12 列六篇專門回應：Devins\ix{persons}{devins@Devins, Neal} \& Fisher\ix{persons}{fisher@Fisher, Louis}、Hartnett\ix{persons}{hartnett@Hartnett, Edward A.}、Nagel\ix{persons}{nagel@Nagel, Robert F.}、Peabody\ix{persons}{peabody@Peabody, Bruce G.}、Sherwin\ix{persons}{sherwin@Sherwin, Emily} 與 Tushnet\ix{persons}{tushnet@Tushnet, Mark} 的 Two Versions of Judicial Supremacy，另附 Harrison\ix{persons}{harrison@Harrison, John} 的未刊稿；註 13 列 Tushnet 的 \en{Taking the Constitution Away from the Courts} 27--29 頁與 Buss、Goldstein、Sherry\ix{persons}{sherry@Sherry, Suzanna} 的文章。批評者視兩人為不顧歷史、為最高法院的傲慢辯護、貶抑國會與行政部門解釋憲法之能力的人。

本文說明目的不在於「critics and \textquotesingle misunderstood\textquotesingle{} authors」〔批評者與「被誤解」的作者〕之間那種針鋒相對的書信往來，兩人以這批批評作為「joining issue」〔就爭點交鋒〕的機會，延續一個涉及憲法的角色與法院在民主社會角色的辯論\pg{458 頁}。

\section{第 I 部分：問題的性質\pg{459 至 467 頁}}
\label{sec:AS-2}

\subsection{（一）先於憲法的憲法\pg{459 至 460 頁，註 14--15}}

第 I 部分 A 節處理兩人所問問題的種類：最高法院應否是憲法的最高解釋者，其解釋應否拘束他人，即使他人自己的解釋看起來更好\pg{459 頁}。他們說這個問題無從由憲法本身回答。Harrison\ix{persons}{harrison@Harrison, John}、Hartnett\ix{persons}{hartnett@Hartnett, Edward A.}、Devins\ix{persons}{devins@Devins, Neal} 與 Fisher\ix{persons}{fisher@Fisher, Louis} 都強調，兩人忽略憲法未指定最高法院為最高解釋者這項具有決定性的資料，也忽略說明並支持此點的憲法史\pg{459 頁，註 14 載三人著作的頁次}。

Alexander\ix{persons}{alexander@Alexander, Larry} 與 Schauer\ix{persons}{schauer@Schauer, Frederick} 為論證之便接受這些歷史證據：制憲者無意使最高法院成為憲法的最高解釋者，各州批准者亦無此理解，文本並未指定。他們再進一步假設憲法實際上否定這樣的角色，但仍主張其論證涉及「this Constitution in this constitutional system」〔這部憲法、這個憲法體制〕，並未停留在假想的體制\pg{459 頁}。理由是他們論證的是一項 preconstitutional\ix{terms}{03@\zhlidxcat{憲法之效力基礎}!先於憲法的憲法} norm〔先於憲法的規範〕，這項規範「determines not what was adopted then, but how what was adopted then should be regarded now」〔決定當時所採行者在今日應被如何看待〕。憲法作為根本法的權威「ultimately rests not on facts about the past, but on the Constitution\textquotesingle s acceptance as authoritative in the present」〔終究取決於它在現在被接受為有權威，與過去的事實無關〕。他們稱此為邏輯上的論點，並說整個論證以它為基礎\pg{459 至 460 頁，註 15 引 Michelman\ix{persons}{michelman@Michelman, Frank I.} 論人民的制憲者身分}。

\subsection{（二）假設的最高法院解釋條款\pg{460 至 461 頁，註 16--18}}

兩人設想一個與事實相反的情形來說明上述論點：憲法第 6 條\ix{statutes}{02@美國!01@聯邦憲法!006@第 6 條}的至上條款不僅寫「This Constitution ... shall be the Supreme Law of the Land」，還明定最高法院對憲法的解釋對其他聯邦與各州法官、行政與立法官員及各州官員「shall be final and binding」〔具有終局拘束力〕\pg{460 頁}。即使憲法文本真是如此，司法至上\ix{terms}{01@\zhlidxcat{司法審查權限之歸屬}!司法至上}的問題仍然未決，因為仍須問誰來解釋該條款，而該解釋是否拘束他人。

這可以用 Wittgenstein\ix{persons}{wittgenstein@Wittgenstein, Ludwig（維根斯坦）} 的說法表達，即規則無法決定其自身的適用（註 16 以箭頭記號為例，問人如何知道該朝箭頭尖端或尾端前進）；也可以用 Kelsen\ix{persons}{kelsen@Kelsen, Hans（凱爾森）} 的 grundnorm\ix{terms}{03@\zhlidxcat{憲法之效力基礎}!基本規範}〔基本規範〕或 Hart\ix{persons}{hart@Hart, H. L. A.（哈特）} 的 ultimate rule of recognition\ix{terms}{03@\zhlidxcat{憲法之效力基礎}!承認規則}〔終極承認規則〕表達\pg{460 至 461 頁，註 17 引 Kelsen「is grounded as valid by the presupposed basic norm」；註 18 引 Hart}。共同的要點是：法律文本的地位無法，也不能在不陷入無窮後退的情形下，由該文本自身決定。判斷文本的地位須離開文本，因此文本也無法回答關於某人對該文本的解釋，包括對其中指定權威解釋者的部分的解釋，是否有權威\pg{461 頁}。

\subsection{（三）「two」與 Conglish\pg{461 至 462 頁，註 19--23}}

兩人承認就這個虛構條款而言，上述論點在許多方面屬形式性的。同一論點也適用於第 3 條第 3 項\ix{statutes}{02@美國!01@聯邦憲法!003@第 3 條}「testimony of two Witnesses」中的 two 是否指 3：在嚴格邏輯上，文本本身無法決定它應被理解為普通英文（two 指 2），或某種 Conglish（two 指 3；註 20 稱此例為 Kripke 對 Wittgenstein\ix{persons}{wittgenstein@Wittgenstein, Ludwig（維根斯坦）} 的討論之變形）\pg{461 頁}。他們接著承認，以普通英文理解憲法是根深蒂固而被廣泛接受的先於憲法\ix{terms}{03@\zhlidxcat{憲法之效力基礎}!先於憲法的憲法}的理解，將 two 解為 3 會違反最穩固的先於憲法的規則之一。

兩人隨即指出，以普通英文解讀文本並不如設想的那樣清楚。實際憲法文本中已有明顯違反文字的解釋：第 11 修正案的「another State」被解釋為包含同一州（註 21 引 Monaco v. Mississippi\ix{cases}{02@美國!01@聯邦最高法院!1934Monaco v. Mississippi（1934）@Monaco v. Mississippi（1934）} 與 Hans v. Louisiana\ix{cases}{02@美國!01@聯邦最高法院!1890Hans v. Louisiana（1890）@Hans v. Louisiana（1890）}），以及依 Stevens 大法官的看法，禁止各州「impairing the Obligation of Contracts」〔損害契約義務〕的規定亦然（註 22 引 Keystone Bituminous Coal Ass\textquotesingle n v. DeBenedictis 的「not to be read literally」）。這些被普遍接受的對文字的違反，被理解為與其餘文本繼續有效並行不悖（註 23 說明無人把它們理解為廢止或取代整部憲法），所以對虛構條款所作的解釋，使其支持的司法解釋至上\ix{terms}{01@\zhlidxcat{司法審查權限之歸屬}!解釋至上}程度低於其字面所示，在邏輯上並非不可能\pg{461 至 462 頁}。由此兩人說：「What counts as law is determined extra-legally, and what counts as the constitution is determined extra-constitutionally」〔何者算是法律，在法律之外決定；何者算是憲法，在憲法之外決定〕，所以偏離文本表面意義者是否有效，無法以循環地援引文本回答\pg{462 頁}。

\subsection{（四）實際憲法的情形與兩種 constitution\pg{462 至 463 頁，註 24--26}}

兩人以 a fortiori〔更何況〕推及實際文本：連如此明確的司法至上\ix{terms}{01@\zhlidxcat{司法審查權限之歸屬}!司法至上}條款都使問題在邏輯上未決，則在一部既未處理司法審查\ix{terms}{01@\zhlidxcat{司法審查權限之歸屬}!司法審查}權本身（註 24 引 Harrison\ix{persons}{harrison@Harrison, John} 與 Thayer\ix{persons}{thayer@Thayer, James Bradley} 1893 年的論文），也未處理司法解釋權威的憲法上，這個問題「remains stunningly open」〔仍然令人驚訝地未決〕\pg{462 頁}。

他們用「preconstitutional\ix{terms}{03@\zhlidxcat{憲法之效力基礎}!先於憲法的憲法}」指邏輯上的先於，與時間上的先後無關\pg{462 至 463 頁，註 25 引 Kay}。憲法文本應如何解讀，在邏輯上先於文本，因為這些問題最終決定文本的意義；因為它們同時構成一個國家與一部憲法，所以先於憲法的問題在某個意義上具有構成性，因而也是憲法性的。要理解大寫 C 的 Constitution，須先理解它所依託的小寫 c 的 constitution\pg{463 頁，註 26 引 Michelman\ix{persons}{michelman@Michelman, Frank I.} 與 Schauer\ix{persons}{schauer@Schauer, Frederick}}。兩人認為誰的解釋應被視為有權威，是這一類先於憲法的問題中最重要的一個。

據此，他們認為 Devins\ix{persons}{devins@Devins, Neal}、Fisher\ix{persons}{fisher@Fisher, Louis}、Hartnett\ix{persons}{hartnett@Hartnett, Edward A.} 與 Harrison 的歷史與文本分析雖具啟發，結論性遠低於這些作者所認為的。兩人以箭頭作類比：若有人問為何朝箭頭尖端開車，回答說箭頭告訴我們朝尖端走，就是沒有理解使箭頭有意義的、箭頭之外的條件。同理，主張應終局地以憲法文本與歷史決定該文本與歷史告訴我們什麼，犯的是同一個概念錯誤\pg{463 頁}。

\subsection{（五）對「既非經驗亦非歷史」一語的修正\pg{464 頁，註 27}}

第 I 部分 B 節承認批評者在先於憲法\ix{terms}{03@\zhlidxcat{憲法之效力基礎}!先於憲法的憲法}的性質上未必同意，更實質的分歧在於用何種分析或證據回答關於解釋權威的問題\pg{464 頁}。排除求諸憲法自身的自我指涉之後，指引在何處？

前文說這項探討「neither empirical nor historical」〔既非經驗的、也非歷史的〕\pg{註 27 引前文 1369 頁}。兩人現在說這樣表述「Perhaps too casually」〔也許過於隨意〕，作為對所作所為的描述（著重制度設計的邏輯與哲學問題）雖然準確，但令許多批評者疑惑：這項探討怎能那樣不涉經驗？兩人的結論是：「the empirical dimension is one that cannot be avoided」〔經驗面向是無法避免的〕，歷史可以是這種經驗探討的構成部分之一，但只是之一。他們說這不改變其主張的性質與所用證據，但對當時的措辭表示後悔，也後悔暗示制度設計的問題可以在經驗真空中回答\pg{464 頁}。

\subsection{（六）歷史與文本的同等地位\pg{464 至 465 頁，註 28--29}}

說探討不涉歷史，指的是：以歷史的權威回答先於憲法\ix{terms}{03@\zhlidxcat{憲法之效力基礎}!先於憲法的憲法}的問題，與依賴憲法文本具有相同地位，終究是 question-begging〔乞題〕\pg{464 頁}。兩人假設發現五十五位制憲會議代表與批准各州的多數議會，都把第 3 條與\ix{statutes}{02@美國!01@聯邦憲法!003@第 3 條}第 6 條理解\ix{statutes}{02@美國!01@聯邦憲法!006@第 6 條}為指涉最高法院的解釋專屬與解釋至上\ix{terms}{01@\zhlidxcat{司法審查權限之歸屬}!解釋至上}，那時可說 Madison\ix{persons}{madison@Madison, James（麥迪遜）} 等人同意了他們的主張，但兩人仍明示不依賴 Madison 的想法決定該問題。

理由仍在自我指涉。註 28 指出這是 Schauer\ix{persons}{schauer@Schauer, Frederick} 先前用過的例子之一的變體：假設兩人寫一部以他們為最高解釋者的憲法，規定簽署並在 Constitutional Commentary 上刊登簽署的通知後生效（註 29 寫著「We signed it.」〔我們簽了〕）。這部憲法自己規定的生效條件已經成就，並依其內部條件有效，但這遠不足以使它成為美國憲法。憲法之為憲法，取決於一個政治共同體在經驗上接受其為憲法；沒有這種「acceptance as social fact」〔作為社會事實\ix{terms}{03@\zhlidxcat{憲法之效力基礎}!社會事實}的接受〕，再多的形式內部效力也不使文件成為憲法\pg{465 頁}。

據此，決定 Madison 的憲法成為美國憲法、Alexander\ix{persons}{alexander@Alexander, Larry} 與 Schauer 的憲法則否的因素，並未含在文件之中，1787 年也無從決定，現在由美國人民決定\pg{465 頁}。1787 年憲法之成為 2000 年憲法，取決於 2000 年的社會事實，則制憲者在 1787 年的意圖於 2000 年具有何種地位，也取決於 2000 年的決定。「The relevance of history is not determined historically, but by a present political and social decision」〔歷史有何相關性，取決於當下的政治與社會決定〕，所以 1787 年先於憲法的問題的判斷，必然是非歷史的問題；這個非歷史的問題完全可能以賦予歷史一定角色的方式回答\pg{465 至 466 頁}。

\subsection{（七）Cooper 規則的來源與歷史的內容價值\pg{466 頁}}

Madison\ix{persons}{madison@Madison, James（麥迪遜）} 對司法至上\ix{terms}{01@\zhlidxcat{司法審查權限之歸屬}!司法至上}的想法，若他真的想過，與文本所說的一樣不具決定性\pg{466 頁}。兩人的論證是：Cooper 規則\ix{terms}{01@\zhlidxcat{司法審查權限之歸屬}!Cooper 規則}在規範上優於他們稱為（兩人承認帶有傾向性）「institutional anarchy」〔制度上的無政府狀態\ix{terms}{02@\zhlidxcat{定紛止爭論}!制度上的無政府狀態}〕的狀態。若原始憲法沒有 Cooper\ix{cases}{02@美國!01@聯邦最高法院!1958Cooper v. Aaron（1958）@Cooper v. Aaron（1958）} 規則，則原始憲法在規範上不如含有該規則的憲法。取得該規則可以經第 5 條的修憲程序\ix{statutes}{02@美國!01@聯邦憲法!005@第 5 條}，兩人主張的途徑則是以取得原始憲法本身的方式取得它，也就是接受它為有權威。他們說：「In Cooper, the Supreme Court ran that rule up the flagpole. If we like it, all we have to do is salute.」〔Cooper 案中，最高法院把這項規則升上旗桿。如果我們喜歡，只須敬禮。〕並說多數人似乎確已敬禮，若然，Cooper 規則就與憲法本身同樣「constitutional」\pg{466 頁}。

退一步問社會是否應該敬禮時，兩人承認先前對歷史的價值否定得過快。歷史不具權威，不等於與權威無關而不具意義。制憲者聰明，與 Plato\ix{persons}{plato@Plato（柏拉圖）}、Aristotle\ix{persons}{aristotle@Aristotle（亞里斯多德）}、Confucius、Hobbes、Locke、Lincoln\ix{persons}{lincoln@Lincoln, Abraham（林肯）}、Stanton、Rawls\ix{persons}{rawls@Rawls, John（羅爾斯）}、Douglass 一樣可以學習；法國、日本、印度、巴西、南非的歷史亦然。兩人說：「what we learn from history we learn because of its content and not because of its authority」〔我們從歷史學到的東西，因其內容而學到，與其權威無關〕。Madison 說過什麼應與探討有關，但離決定問題還差得遠\pg{466 頁}。

\subsection{（八）規範性的制度設計與經驗證據\pg{466 至 467 頁}}

先於憲法\ix{terms}{03@\zhlidxcat{憲法之效力基礎}!先於憲法的憲法}的問題是規範性的制度設計問題，亦即應採行何種制度設計，不能與世界上的事實無關地探討。兩人因此說，主張探討不涉經驗「was also then and is now a mistake」〔當時與現在都是錯誤〕。正確的主張是：問題的地位，即其先於憲法的性質，本身是邏輯問題，與歷史或經驗無關；問題的地位確立之後，要判斷某種制度設計是否較佳時，無視來自外部世界的證據，包括歷史證據，是愚蠢的\pg{466 至 467 頁}。他們認為，推論某些誘因與制度設計會導致某些行為屬於適當，但這些是經驗推論，應以實際行為的證據補充與檢驗，證據可以是歷史的、社會學的、統計的、心理學的或其他性質\pg{467 頁}。

\section{第 II 部分：定紛止爭功能及其最適所在\pg{467 至 478 頁}}
\label{sec:AS-3}

\subsection{（一）定紛止爭是否被高估\pg{467 至 468 頁，註 30}}

第 II 部分 A 節重述核心主張：法律，尤其憲法，主要功能之一是為了定紛止爭\ix{terms}{02@\zhlidxcat{定紛止爭論}!定紛止爭}而就應為之事提供某種程度的協調性定紛。在道德與政治歧見並存的世界，法律可以提供一種「neither final nor conclusive, but nevertheless authoritative」〔既非終局、也非結論性，但仍有權威〕的定紛，對一階\ix{terms}{02@\zhlidxcat{定紛止爭論}!一階定紛與二階定紛}歧見提供二階的解決，使決定得以作成、行動得以協調\pg{467 頁}。

批評者分兩種。第一種以 Tushnet\ix{persons}{tushnet@Tushnet, Mark} 為最，其他人大多也屬之，在他們背後還有龐大的傳統，在公共的憲法論述中看到獨立的價值，認為這種論述不經持續的公共爭議與辯論就不會出現\pg{467 至 468 頁}。註 30 說明這個傳統涉及的文獻可舉出數以千計，並列 Ackerman \en{We the People}、Levinson\ix{persons}{levinson@Levinson, Sanford} \en{Constitutional Faith}、Sunstein \en{The Partial Constitution}、Lipkin、Michelman\ix{persons}{michelman@Michelman, Frank I.} 與 Minow。依此類主張，憲法的首要價值在於它是公共審議的場域與焦點，使犯罪控制、平權措施、性別平等、性傾向、權力下放給州、總統行為、宗教在公共生活中的角色等議題得到辯論；過度強調憲法的定紛功能，會截斷憲法傳統核心的公共辯論與審議\pg{468 頁}。

\subsection{（二）公共討論與成文憲法的關係\pg{468 頁，註 31--33}}

兩人的回應是「We demur」〔我們不同意〕。他們不否認公共論述的價值（註 31 補「At least here」，並指 Schauer\ix{persons}{schauer@Schauer, Frederick} 另有 Discourse and Its Discontents 討論其局限），但認為以為美國人需要憲法才能辯論平權措施、刑事司法、墮胎、宗教與國家、隱私、死刑，是美國憲法學者的 conceit〔自負〕。這些辯論在沒有單一成文憲法文件的國家，主要是英國、紐西蘭與以色列，同樣興盛（註 32 引 Finer 等與 Jackson \& Tushnet\ix{persons}{tushnet@Tushnet, Mark} 的憲法比較著作）；在成文憲法並未將這些特定議題憲法化的國家亦然（註 33 引 Bell 論法國憲法）。他們承認要權威地回答這個問題，須有以憲政各面向為變項的比較審議研究，而這種研究尚不存在，但依其看法，憲法是否為對重大政治與道德問題的豐富公共辯論的必要條件、充分條件，甚至顯著的因果貢獻者，「far from certain」〔遠不確定〕\pg{468 頁}。

\subsection{（三）法律的特殊功能\pg{468 至 469 頁，註 34}}

兩人說明其論證的前提：同意優於歧見、定紛\ix{terms}{02@\zhlidxcat{定紛止爭論}!定紛止爭}優於保持開放，雖然兩種世界觀都有道理，卻都不在其內；其前提是法律所履行的特殊功能，即同意與定紛似乎是法律的特有領域，也是被以近似法律的方式理解的成文憲法的特有領域\pg{469 頁}。另有以討論為焦點的公共文件，例如孔子或毛主席的著作，兩人的論證以美國憲法與這類文件之間，以及憲法與 Declaration of Independence 之間有重要差異為前提\pg{469 頁}。

註 34 指出，Tushnet\ix{persons}{tushnet@Tushnet, Mark} 在 \en{Taking the Constitution Away from the Courts} 中特別以憲法序言與獨立宣言這類非法律文件作為公共討論基本價值的焦點。兩人的問題是把重大事項的公共辯論與憲政主義兩個問題混為一談是否明智：若兩者相同，Tushnet 對而兩人錯；若不同，差別必在公共論述之外，即對該論述或其產物的限制之中。

\subsection{（四）憲法文本中已定紛的事項與 Brandeis 之語\pg{469 至 470 頁，註 35--37}}

憲法處理的事項通常比靠左或靠右行車更具道德分量，但為定紛\ix{terms}{02@\zhlidxcat{定紛止爭論}!定紛止爭}而定紛的價值並不因此在憲法中缺席。兩人列舉人們可能有嚴肅道德與政治歧見的事項：比例代表制或「first past the post」〔單一選區相對多數決〕、有無國教、民事或刑事案件是否採陪審團、刑事被告能否被迫作證（違者以藐視法庭處罰或可推論有罪）、總統可連任幾個四年任期。這些在美國憲法文本中大致已有定紛，從議程上移除，連不同意其實質內容的人也一樣\pg{469 頁}。法律被當作有權威時，為不同意定紛內容的人提供服從的 content-independent\ix{terms}{02@\zhlidxcat{定紛止爭論}!與內容無關之理由} reason\pg{469 至 470 頁}。

兩人引 Brandeis\ix{persons}{brandeis@Brandeis, Louis D.（布蘭戴斯）} 大法官在 Burnet v. Coronado Oil\ix{cases}{02@美國!01@聯邦最高法院!1932Burnet v. Coronado Oil \& Gas Co.（1932）@Burnet v. Coronado Oil \& Gas Co.（1932）} \& Gas Co.（1932）不同意見\ix{terms}{04@\zhlidxcat{法院之功能}!不同意見與協同意見}書的話：「in most matters it is more important that the applicable rule of law be settled than that it be settled right」〔多數事項上，適用的法律規則被確定，比被確定得正確更為重要〕\pg{470 頁，註 36}。他們指出，有人主張強力的先例拘束在憲法脈絡中重要性較低（註 37 引 Webster 與 Payne 兩案及 Cooper\ix{cases}{02@美國!01@聯邦最高法院!1958Cooper v. Aaron（1958）@Cooper v. Aaron（1958）}），但憲法文本本身強力地支持 Brandeis 的觀點，制憲者認為為定紛而定紛的約束，往往比面對世界變遷時更大的彈性更有價值\pg{470 頁}。

註 35 回應 Hartnett\ix{persons}{hartnett@Hartnett, Edward A.}（並言及 Nagel\ix{persons}{nagel@Nagel, Robert F.}）認為兩人反對歧見、勸人服從即是譴責歧見的主張，說不知如何理解；承認前文用語「without regard」有歧義，可能促成誤解，但以每小時 55 英里的速度行駛並公開批評任何低於 70 英里的限速，兩者可以並行，同樣，對最高法院解釋的公開異議，與服從該解釋直到法院改變看法，完全相容。

\subsection{（五）文本未定紛的爭議與最高法院意見\pg{470 至 472 頁，註 38--47}}

上述例子是 biased sample〔有偏的樣本〕，因為都是文本已相當明確的情形。美國憲法文本如此不確定，更常見的是帶有憲法面向而文本本身未提供定紛\ix{terms}{02@\zhlidxcat{定紛止爭論}!定紛止爭}的爭議\pg{470 頁}。兩人列了一串問題：

\begin{enumerate}
\def\labelenumi{\arabic{enumi}.}
\tightlist
\item
  增修條文第 1 條\ix{statutes}{02@美國!01@聯邦憲法!101@增修條文第 1 條}的禁止設立國教條款是否禁止公立學校中由教師組織的、無教派色彩的禱告（註 38 引 Wallace v. Jaffree\ix{cases}{02@美國!01@聯邦最高法院!1985Wallace v. Jaffree（1985）@Wallace v. Jaffree（1985）}、Abington Sch. Dist. v. Schempp\ix{cases}{02@美國!01@聯邦最高法院!1963Abington School District v. Schempp（1963）@Abington School District v. Schempp（1963）}、Engel v. Vitale\ix{cases}{02@美國!01@聯邦最高法院!1962Engel v. Vitale（1962）@Engel v. Vitale（1962）}），或畢業典禮上由學生主導者（註 39 引 Lee v. Weisman\ix{cases}{02@美國!01@聯邦最高法院!1992Lee v. Weisman（1992）@Lee v. Weisman（1992）}）。
\item
  平等保護條款是否禁止各州維持單一性別的大學（註 40 引 Mississippi Univ. for Women v. Hogan\ix{cases}{02@美國!01@聯邦最高法院!1982Mississippi University for Women v. Hogan（1982）@Mississippi University for Women v. Hogan（1982）} 與 United States v. Virginia\ix{cases}{02@美國!01@聯邦最高法院!1996United States v. Virginia（1996）@United States v. Virginia（1996）}）。
\item
  第 4 修正案的令狀要求是否禁止在審判中使用未經令狀取得的、具證明力的證據（註 41 引 Mapp v. Ohio\ix{cases}{02@美國!01@聯邦最高法院!1961Mapp v. Ohio（1961）@Mapp v. Ohio（1961）}）。
\item
  言論自由條款是否涵蓋明示鼓勵種族暴力的言論（註 42 引 Brandenburg v. Ohio\ix{cases}{02@美國!01@聯邦最高法院!1969Brandenburg v. Ohio（1969）@Brandenburg v. Ohio（1969）}）。
\item
  死刑是否殘酷而不尋常（註 43 引 Gregg v. Georgia\ix{cases}{02@美國!01@聯邦最高法院!1976Gregg v. Georgia（1976）@Gregg v. Georgia（1976）}）。
\item
  國會未明示行動時，各州能否對州際商業課以財務或後勤上的負擔（註 44 引 Kassel）。
\item
  現代立法實務的需要是否容許國會將立法職責委由行政機關行使（註 45 引 Mistretta 與 American Textile），或省略兩院通過與送交總統的正式程序（註 46 引 INS v. Chadha\ix{cases}{02@美國!01@聯邦最高法院!1983INS v. Chadha（1983）@INS v. Chadha（1983）}）。
\end{enumerate}

這些場合都有合理的道德與政治歧見而無清楚的文本定紛，取而代之的是最高法院試圖解決該議題的意見，其解決當然不消除根本的歧見\pg{471 頁}。兩人的核心主張是：在這種情形，為定紛而定紛的價值，使最高法院以外的機構，尤其是下級法院、州議會、國會與總統，應把該解決視為有權威，即使它們繼續不同意其實質內容。承認最高法院的權威與解釋至上\ix{terms}{01@\zhlidxcat{司法審查權限之歸屬}!解釋至上}，這些機構就如同承認憲法文本本身的權威與至上，對穩定與社會和諧作出貢獻\pg{471 頁}。

兩人以 Clinton 總統為例：若他承認服從第 22 修正案、不競選第三任期具有定紛的價值，則同理他也應承認拒絕簽署一部明顯違反最高法院近期一致見解的 Communications Decency Act 具有定紛的價值\pg{471 頁，註 47 對照 Reno v. ACLU\ix{cases}{02@美國!01@聯邦最高法院!1997Reno v. ACLU（1997）@Reno v. ACLU（1997）} 與 Sable Communications v. FCC\ix{cases}{02@美國!01@聯邦最高法院!1989Sable Communications v. FCC（1989）@Sable Communications v. FCC（1989）}}。他們進一步說：若定紛論證失敗，則憲法的大部分難以解釋；若成功，則它支持單一權威解釋者，其強度與它支持單一權威憲法相同\pg{471 至 472 頁}。

\subsection{（六）Tushnet、Peabody 與 Sherwin 的批評\pg{472 頁，註 48--50}}

兩人說 Tushnet\ix{persons}{tushnet@Tushnet, Mark} 對法律式定紛\ix{terms}{02@\zhlidxcat{定紛止爭論}!定紛止爭}作為超越性價值採取最強烈的反對立場\pg{472 頁，註 48 引 Tushnet, Two Versions of Judicial Supremacy 952--57 頁}；Peabody\ix{persons}{peabody@Peabody, Bruce G.} 同樣批評兩人過度強調定紛（註 49），雖承認其部分保留，仍主張兩人對定紛與穩定的推崇貶低了其他重要價值，包括實體與程序的，如創造性及迅速適應變動世界所需的彈性\pg{472 頁，註 50}。

註 50 另處理 Peabody 與 Sherwin\ix{persons}{sherwin@Sherwin, Emily} 的批評。Sherwin 對 Cooper\ix{cases}{02@美國!01@聯邦最高法院!1958Cooper v. Aaron（1958）@Cooper v. Aaron（1958）} 的辯護沒有異議，她的批評在於兩人如何調和該辯護與他們對 Lincoln\ix{persons}{lincoln@Lincoln, Abraham（林肯）} 有義務服從 Dred Scott\ix{cases}{02@美國!01@聯邦最高法院!1857Dred Scott v. Sandford（1857）@Dred Scott v. Sandford（1857）} 案命令所說的話。兩人主張 Lincoln 有理由，或至少可以有理由不服從法院，這項讓步不損及對 Cooper 的辯護；Sherwin 則認為那樣兩人所主張的實際上只是服從的一項考慮因素，算不上真正的規則，而唯有真正的規則才能提供穩定、可預測與協調的定紛利益。兩人回應，同意強力規則的必要性，但認為她低估了權威的不對稱\ix{terms}{02@\zhlidxcat{定紛止爭論}!權威的不對稱}：權威的規則主張絕對的服從，不是僅供考慮，要求受拘束者即使確信衡量各項理由後，包括有規則及其定紛利益的理由，不服從為宜，仍須服從；受拘束者面對權威規則時，卻不能讓渡自己的理性與道德主體性。因此權威所應命令的與主體所應為的之間，永遠有差距，說法院從其立場有理由要求 Lincoln 服從，而 Lincoln 從其立場在道德上有理由不服從，並無不一致\pg{472 頁，註 50}。

\subsection{（七）法律的職責與憲政主義的價值\pg{472 至 473 頁}}

兩人說自己的論證隱含一點，現在要明示：法律、法院或憲法不是，也不應是制度設計中一切善的寄託。他們以汽車煞車為喻：設計煞車時不應以速度、加速度或操控的容易為目標，即使承認這些目標也相關。同理，兩人依賴的前提是法律性質的制度比其他社會制度更能發揮促進穩定與定紛\ix{terms}{02@\zhlidxcat{定紛止爭論}!定紛止爭}的功能，在複雜而分化的社會中這不一定是錯誤的職能分工觀念\pg{472 至 473 頁}。主張定紛是法律的特殊職責，並不是主張所有社會制度有同樣職責，或定紛在社會價值中的地位，應與它在法律這個社會子集中同樣重要\pg{473 頁}。

兩人據此說，從定紛來看司法至上\ix{terms}{01@\zhlidxcat{司法審查權限之歸屬}!司法至上}，就是把憲法視為提供穩定與約束的功能。若穩定與約束在憲法價值中不居首位，他們不明白為何一開始要對一階\ix{terms}{02@\zhlidxcat{定紛止爭論}!一階定紛與二階定紛}政治決定施加二階的憲法約束。若穩定與約束是說明憲政主義本身的核心，則主張單一權威解釋者的論證「flows directly from the deepest values of constitutionalism」〔直接源於憲政主義最深層的價值〕\pg{473 頁}。

\subsection{（八）法院是否為最佳的定紛者\pg{473 至 474 頁，註 51}}

第 II 部分 B 節處理的批評，把除 Sherwin\ix{persons}{sherwin@Sherwin, Emily} 之外所有批評者連在一起：兩人從定紛\ix{terms}{02@\zhlidxcat{定紛止爭論}!定紛止爭}本身是好事的前提，過快地推到法院，尤其是最高法院，是履行此功能的最佳機構。這些批評者承認定紛是好事，甚至承認它可能需要單一權威解釋者，但質疑為何應交給最高法院，而非其他機構\pg{473 頁}。

論證之一強調最高法院與國會、與行政部門之間的對話。Devins\ix{persons}{devins@Devins, Neal} 與 Fisher\ix{persons}{fisher@Fisher, Louis} 為兩人稱為「institutional anarchy」的狀態的暫時形式辯護，主張法院與政治部門之間比 Cooper\ix{cases}{02@美國!01@聯邦最高法院!1958Cooper v. Aaron（1958）@Cooper v. Aaron（1958）} 所容許的複雜得多的往返是好事，長期會產生較好的憲法決定，因而帶來政治穩定的提高\pg{473 頁}。兩人說自己沒有做過回應這項命題所需的完整社會科學檢驗，因此無法直接反駁。若 Devins 與 Fisher 的意思是所有權威都脆弱地建立在受拘束者的接受之上，權威所提供的 content-independent\ix{terms}{02@\zhlidxcat{定紛止爭論}!與內容無關之理由} reasons 因深刻的弔詭而同樣脆弱，兩人同意\pg{473 至 474 頁；註 51 說明 Alexander\ix{persons}{alexander@Alexander, Larry} 在多篇文章中多次強調此點，列 Can Law Survive the Asymmetry of Authority? 等}。

以下（九）至（十三）在刊本 474 至 478 頁，與下一節丁\pg{478 至 481 頁}同屬課綱指定頁次\pg{455 至 473 頁、481 至 482 頁}以外的部分。

\subsection{（九）公眾對權威定紛的接受\pg{474 頁，註 52--54}}

若 Devins\ix{persons}{devins@Devins, Neal} 與 Fisher\ix{persons}{fisher@Fisher, Louis} 的意思是不應權威地定紛\ix{terms}{02@\zhlidxcat{定紛止爭論}!定紛止爭}憲法問題，因為這樣做冒著損害權威所賴以存在的接受之風險，兩人則表示反對。公眾可以也常常接受權威的定紛，包括憲法上的定紛，即使在道德、政策與文本解釋上強烈不同意\pg{474 頁}。兩人舉例：早期 New Deal 階段最高法院宣告為對抗大蕭條而制定的聯邦立法無效的裁判（註 52 列 United States v. Butler\ix{cases}{02@美國!01@聯邦最高法院!1936United States v. Butler（1936）@United States v. Butler（1936）} 與 A.L.A. Schechter Poultry\ix{cases}{02@美國!01@聯邦最高法院!1935A.L.A. Schechter Poultry Corp. v. United States（1935）@A.L.A. Schechter Poultry Corp. v. United States（1935）} Corp. v. United States），以及 Warren 法院宣告學校禱告違憲的裁判（註 53 指向註 39），在實質上都極不受歡迎，在解釋上也被廣泛批評為錯誤，但都沒有被廣泛認為缺乏權威。

兩人說權威可能只建立在受拘束者逐時逐刻的接受之上，但這種接受可以承受相當程度的、對政策與解釋的深刻歧見。又，若對司法角色的接受與對司法的尊重\ix{terms}{05@\zhlidxcat{解釋理論與自制}!尊重}主要取決於裁判的實質內容（註 54 引 Adamany \& Grossman、Caldeira 與 Gibson 等人的研究），則法院角色結構上的差異，對其被接受的程度與正當性可能影響不大\pg{474 頁}。

\subsection{（十）解釋分歧與實體結果分歧的區分\pg{474 至 475 頁，註 55}}

兩人主張須區分憲法歧見的種類。分歧僅在於最高法院對憲法規則的解釋，不涉實體結果，也就是人們認為最高法院的憲法規則優於制憲者採用的規則時，Cooper 規則\ix{terms}{01@\zhlidxcat{司法審查權限之歸屬}!Cooper 規則}的理由最強\pg{474 至 475 頁，註 55 補充不宜高估人民對最高法院判決所知所信}。這不是說法院不應忠實解釋憲法，即使法院認為憲法在道德或政策上不是最佳甚至有害；僅指出利害只在於對先前定紛\ix{terms}{02@\zhlidxcat{定紛止爭論}!定紛止爭}的解釋是否正確時，定紛的價值似乎重於解釋正確的價值\pg{475 頁}。

政治部門認為最高法院的解釋不僅不正確，而且不公正或政策不良時，定紛是否應讓位於實體正確？兩人從制度設計的立場，也就是貫穿其整個計畫的立場，仍答「no」\pg{475 頁}。理由是許多，或許大多數憲法議題涉及實質重要的事項，所以幾乎所有對法院的解釋性歧見都伴隨政策歧見。沒有權威的定紛，大部分憲法議題就維持開放：任何不同意最高法院憲法裁判的政府行動者，都會覺得在法律與道德上，即使不一定在政治或審慎\ix{terms}{05@\zhlidxcat{解釋理論與自制}!審慎}考量上，可以不理會該裁判。兩人舉例：不同意第 4 修正案令狀法現狀的警察，若認為不聲請令狀也能達成目的，就會略過；下級法院法官若無遵守最高法院裁判的法律義務，他達成目的的機會更高；州議會想要公立學校的有組織禱告、市政府想規制猥褻之外的不雅內容、行政官員想運用權力勸阻婦女墮胎，情形亦然\pg{475 頁}。

\subsection{（十一）批評者未主張的範圍\pg{475 至 476 頁，註 56--58}}

兩人說批評者沒有人主張下級法院無須遵循最高法院（註 56 保留 Hartnett\ix{persons}{hartnett@Hartnett, Edward A.} 為可能的例外，指向第 III 部分；註 57 引 Rogers 與 Fried），也無人主張 Meese\ix{persons}{meese@Meese, Edwin, III} 司法部長十二年前主張的立場，即各州不論在行政、立法、執法或司法的身分，均得把最高法院裁判視為只限於該案、不為其他行動者設定法律\pg{475 至 476 頁，註 58 引 Meese 與 Tulane Law Review 專題}。批評者不願強烈支持把解釋權分散給 50 個州，顯示辯論的範圍比乍看窄，兩人對國家憲法的全國解釋至上\ix{terms}{01@\zhlidxcat{司法審查權限之歸屬}!解釋至上}的論證已為人廣泛接受\pg{476 頁}。

仍然，Devins\ix{persons}{devins@Devins, Neal} 與 Fisher\ix{persons}{fisher@Fisher, Louis}、Nagel\ix{persons}{nagel@Nagel, Robert F.}、Peabody\ix{persons}{peabody@Peabody, Bruce G.} 與 Tushnet\ix{persons}{tushnet@Tushnet, Mark} 都主張，需要權威定紛\ix{terms}{02@\zhlidxcat{定紛止爭論}!定紛止爭}並不決定哪個機構應擁有該權威。民主組成的機構也能定紛憲法議題，所以對某些或許多議題，國會在履行定紛功能上可能優於法院\pg{476 頁}。

\subsection{（十二）最高法院作為定紛者的理由\pg{476 至 478 頁，註 59--61}}

兩人提出的理由之一，是法院「relative insulation from political winds」〔與政治風向的相對隔絕〕，除非認為憲法解釋就是且應是當代價值與政策的表達，否則這是明顯的優點，使最高法院優於國會或行政部門\pg{476 頁}。Nagel\ix{persons}{nagel@Nagel, Robert F.} 特別指責兩人誇大最高法院的穩定程度\pg{註 59 引 Nagel 857--60 頁}，兩人回應，穩定與否只能以兼顧變動頻率與變動幅度的函數判斷。他們舉住宅電流為例：在 108 至 112 伏之間不斷波動的電流，在多數用途上比七天中有六天在 109.9 至 110.1 伏、一天在 62 或 154 伏的電流更穩定。因此只看最高法院不穩定的次數而不比較其與其他部門的變動幅度，如 Nagel 似乎做的，是錯誤的；最高法院的不穩定若頻繁但幅度窄（註 60 引 Blasi 編的 \en{The Burger Court: The Counter-Revolution That Wasn\textquotesingle t}），就更是如此\pg{476 頁}。

兩人的懷疑（承認尚缺他們與批評者都未做過的受控研究）是：最高法院對先例的遵循，使它能比沒有類似先例拘束實踐的政治部門更持久地定紛\ix{terms}{02@\zhlidxcat{定紛止爭論}!定紛止爭}憲法議題\pg{476 至 477 頁，註 61 引 Scalia、Tribe\ix{persons}{tribe@Tribe, Laurence H.} 與 Monaghan}。即使最佳的制度設計是免除法院遵循自己先例的義務（註 62，見下一小節），法院仍優於政治部門：憲法議題經由已建立且有約束力的程序進入法院，政治部門沒有可比的程序；法院作成裁判時附具意見，政治部門不附（雖然可以）；法院只有九位成員且終身任職；法院不能選擇自己的議程，國會可以\pg{477 頁}。這些與其他特徵是否使法院在最終經驗分析上比國會或行政部門更穩定，「is more of an empirical question than we first acknowledged」〔比我們起初承認的更是經驗問題〕，但在缺乏接近解決此問題的非軼事性證據之前，兩人仍不相信國會能比最高法院提供更多為定紛而定紛的定紛\pg{477 至 478 頁}。

\subsection{（十三）最高法院本身是否受 Cooper 規則拘束\pg{477 頁，註 62}}

註 62 提出問題：最高法院自己是否受 Cooper 規則\ix{terms}{01@\zhlidxcat{司法審查權限之歸屬}!Cooper 規則}拘束，也就是它現在認為錯誤的裁判是否仍永久地定紛\ix{terms}{02@\zhlidxcat{定紛止爭論}!定紛止爭}該議題，水平先例的強度應有多大。註中列三種立場：

\begin{enumerate}
\def\labelenumi{\arabic{enumi}.}
\tightlist
\item
  Cooper\ix{cases}{02@美國!01@聯邦最高法院!1958Cooper v. Aaron（1958）@Cooper v. Aaron（1958）} 規則完全適用於法院自己。法院一旦裁判憲法事項，即永久確定；不論法院現在認為先前的裁判多麼錯誤、有害或不公正，先前裁判對法院和所有其他官員都與憲法一樣具有權威。
\item
  法院是唯一免於 Cooper 規則的機構；認為先前憲法先例錯誤時，可以自由推翻。
\item
  法院僅在先例既是錯誤的憲法解釋、又經法院認為不公正或有害時，才可推翻；憲法錯誤是必要但不充分的條件。
\end{enumerate}

兩人評論第一種立場將定紛置於一切對立價值之上，法院只有一次正確解釋特定憲法條文的機會，補救最高法院錯誤的唯一方式是漫長而困難的修憲\ix{terms}{03@\zhlidxcat{憲法之效力基礎}!修憲}。第二種則不夠重視定紛。兩人傾向第三種的某個變體，因為它給法院一道對有害憲法錯誤的安全閥。第三種也有缺點：法院須作道德與政策上的一階\ix{terms}{02@\zhlidxcat{定紛止爭論}!一階定紛與二階定紛}判斷，並犧牲部分定紛價值、容許部分憲法錯誤。故兩人對第三種立場的支持「less than wholehearted and quite tentative」〔不全心，且頗為保留〕\pg{477 頁，註 62}。

\section{第 III 部分：最高法院的功能\pg{478 至 481 頁}}
\label{sec:AS-4}

\subsection{（一）案件裁判者的圖像與 Hartnett\pg{478 頁，註 63}}

第 III 部分指出，許多批評的基底，是以案件裁判者\ix{terms}{04@\zhlidxcat{法院之功能}!案件裁判者與造法者}看待最高法院的圖像，造法者的身分不在其內\pg{478 頁}。兩人說這個圖像符合傳統的法院功能觀，不因此必然無效；它在 Hartnett\ix{persons}{hartnett@Hartnett, Edward A.} 的批評中最明確，Hartnett 認為法院所做的是裁判具體案件，其餘包括意見在內都是「gravy」〔佐料〕。對 Hartnett，司法功能的核心是裁判案件，提出意見居於其外；兩人的整個立場既然建立在意見可被視為提供某種指引之上，在 Hartnett 看來便是誤解了裁判的性質。Hartnett 並以令人印象深刻的歷史論證支持，主張最高法院自建國之初就被構想為案件裁判者，沒有被構想為造法者或意見提供者\pg{478 頁，註 63 列 Hartnett 126--36、146--58 頁，並指 Devins\ix{persons}{devins@Devins, Neal} \& Fisher\ix{persons}{fisher@Fisher, Louis} 91 頁有簡短相似主張，另列 Eisgruber、Merrill 與 Wechsler\ix{persons}{wechsler@Wechsler, Herbert}}。

\subsection{（二）每年受理案件數與法院的角色\pg{478 至 479 頁，註 64--65}}

兩人不爭論 Hartnett\ix{persons}{hartnett@Hartnett, Edward A.} 的歷史，但與對待 Devins\ix{persons}{devins@Devins, Neal} 和 Fisher\ix{persons}{fisher@Fisher, Louis} 的歷史一樣，對該歷史的地位有不同看法。把歷史視為有趣但非決定性的，則可問何種法院應做什麼的圖像，最符合兩人的定紛\ix{terms}{02@\zhlidxcat{定紛止爭論}!定紛止爭}功能觀\pg{478 頁}。他們說去年最高法院依 certiorari 程序以其自己對慎重審酌的定義，慎重審酌的案件只有 94 件，被請求慎重審酌而選擇不審的有 7015 件（註 64 引 \en{The Supreme Court, 1998 Term---The Statistics}, 113 Harv. L. Rev. 400, 407）。若最高法院主要是案件裁判者\ix{terms}{04@\zhlidxcat{法院之功能}!案件裁判者與造法者}，意見提供者至多居次，它「appears to have achieved a level of inefficiency unmatched in any other government agency we know of」〔似乎達到了我們所知任何其他政府機構都比不上的低效率〕。若法院只處理 1\% 的案件量，以求在這 1\% 上作出好的裁判，並附上依 Hartnett 等人看來不太重要的意見，則法院的角色與運作方式「the time has come for major reconsideration」〔該到全面重新考慮的時候了〕\pg{478 至 479 頁}。

註 65 引 Estreicher 與 Sexton 的 \en{A Managerial Theory of the Supreme Court\textquotesingle s Responsibilities} 主張最高法院應更把自己視為下級法院體系的管理者，兩人不反對，另提出更激進的提議：法院也應把自己視為「manager of constitutional signals」〔憲法訊號的管理者〕與「manager of a system of constitutional compliance」〔憲法遵從體系的管理者〕。

\subsection{（三）意見的數量與品質\pg{479 頁，註 66--68}}

兩人說重新考慮或許適當，但原因在於法院寫的意見不夠，或寫得不夠好，與寫了太多意見無關\pg{479 頁，註 66 引 Easterbrook 與 Schauer\ix{persons}{schauer@Schauer, Frederick} 的 Opinions as Rules}。法院作為憲法解釋者是兩人的圖像，與 Hartnett\ix{persons}{hartnett@Hartnett, Edward A.} 的法院作為爭端解決者\ix{terms}{04@\zhlidxcat{法院之功能}!爭端解決者}的圖像相對；前者看到法院需要比目前意見的數量與風格所顯示的，更關心指導、引導、協助以及命令其他機構與部門（註 67 稱此與 Post \& Salop 的看法大致一致）。

如果這是或應當是法院關心的事，兩人預期會看到更清楚的規則、更少沒有多數意見的分歧判決、大法官更關心法院以單一聲音發言，甚至勝過強調自己的論點或堅持自己對結果的看法，更重視 stare decisis，總之更多符合法院作為造法者\ix{terms}{04@\zhlidxcat{法院之功能}!案件裁判者與造法者}的行為，而法院「undeniably is, and in our judgment inevitably must be」〔無可否認就是，且依我們的判斷無可避免必然是〕造法者\pg{479 頁，註 68 引 Rogers，並指 Sherry\ix{persons}{sherry@Sherry, Suzanna} 同意兩人對 Cooper\ix{cases}{02@美國!01@聯邦最高法院!1958Cooper v. Aaron（1958）@Cooper v. Aaron（1958）} 的立場，並認為兩人的立場意味著個別大法官在許多脈絡，例如 City of Boerne\ix{cases}{02@美國!01@聯邦最高法院!1997City of Boerne v. Flores（1997）@City of Boerne v. Flores（1997）}，應如同他們期待其他官員那樣受法院裁判拘束}。

\subsection{（四）Hartnett 論證的最強處\pg{479 至 480 頁，註 69--70}}

Hartnett\ix{persons}{hartnett@Hartnett, Edward A.} 立場最強的論證，在兩人看來出自法院目前作為規則制定者的表現，他的歷史分析雖然正確，但大致無關\pg{479 至 480 頁}。若對權威規則的權威解釋是規則制定的功能，兩人相信如此，則法院的許多運作方式需要改革。兩人說，Hartnett 最有說服力的主張，即使不是他明白提出的，是：依他的圖像法院可以照現在這樣繼續，依兩人的圖像則須花更多時間思考如何對國會、行政部門、下級法院與各州給予更好的指引\pg{480 頁}。

註 69 說明兩人不接受 Hartnett 認為法院「have little business decreeing」〔不太有權宣示〕規則，也不同意意見，也就是理由，無法被遵循。理由作為對某類案件（包含眼前這件）應如何處理的陳述，本質上就是規則，可以如規則般被遵循；例如「I hold in your favor because people have a right to protest in public parks」〔我判你勝訴，因為人們有權在公園抗議〕這句話所含的理由，即構成一條公園內抗議必須獲准的規則。在兩人看來，法院每提出一個判決的理由就是制定規則，所以每寫一份意見都是；主張法院應退出這門事業才是真正激進的主張，他們的建議較溫和，只是請法院承認它不可避免地從事的規則制定，作為更好地做它現在做、且已做了數個世紀之事的第一步。

註 70 補充：若從頭開始，這也可能是要求一部不那麼晦澀的憲法文本的理由。南非 114 頁的新憲法與巴西 211 頁的新憲法也許是極端例子，但當今世界所有制憲過程中，沒有國家選擇像美國憲法那樣簡短晦澀的成文憲法。部分原因可能是美國憲法在祕密中制定，而多數現代憲法，尤其是這兩部，是在非常公開的情形下制定，政治誘因傾向全面納入各種權利、利益與價值；部分也可能是人們愈益認識到，憲法既要作為法律，便不足為奇會看起來像法律。

\subsection{（五）最高法院自身的宣稱\pg{480 至 481 頁，註 71--72}}

兩人把最高法院自己宣稱的解釋至上\ix{terms}{01@\zhlidxcat{司法審查權限之歸屬}!解釋至上}，視為與原始意圖同樣無關，因為這種自我提拔式的論證無法回答其他部門是否應尊重\ix{terms}{05@\zhlidxcat{解釋理論與自制}!尊重}該宣稱，註 71 說明此外並無來自司法以外的政治誘因支撐，例如不理會即使不受歡迎的最高法院裁判的民選官員所要承擔的政治責任，在現行政治環境中顯然不存在\pg{480 頁}。然而另一意義上，法院自己的宣稱高度相關：如果法院自稱是其他官員與部門的引導者，Cooper\ix{cases}{02@美國!01@聯邦最高法院!1958Cooper v. Aaron（1958）@Cooper v. Aaron（1958）} 與 City of Boerne\ix{cases}{02@美國!01@聯邦最高法院!1997City of Boerne v. Flores（1997）@City of Boerne v. Flores（1997）} 的陳述似乎如此主張，則法院的表現可以適當地以其自己宣布的功能衡量\pg{480 至 481 頁，註 72 引 Krotoszynski 論法院如何更有效地履行此功能}。依兩人所支持的功能定義，法院「plainly has a long way to go」〔顯然還有很長的路要走〕\pg{481 頁}。

\subsection{（六）諮詢意見禁令與裁判理由\pg{481 頁，註 73--74}}

兩人預期想像這個角色似乎與禁止諮詢意見\ix{terms}{04@\zhlidxcat{法院之功能}!諮詢意見}相抵觸（註 73 列 Dorf、Field、Frankfurter 與 Clovis \& Updegraff）。他們答覆：若諮詢意見禁令被理解為禁止對於案件結果不必要的、關於法律的有用陳述，法院將只剩宣告判決，或許，僅僅是或許，附上案件事實的陳述。法院一旦超出此範圍，不論稱之為提出意見、說明理由或附帶意見，都在從事提供諮詢意見的事業。禁止法院在個案或爭議之外這樣做，可能有良好的第 3 條理由\ix{statutes}{02@美國!01@聯邦憲法!003@第 3 條}（註 74 引 Krotoszynski 與 Lee 關於 mootness 的文章），但個案或爭議一旦存在，要消除法院所從事的全部說理實踐已為時過晚，這些實踐必然「decide」眼前案件以外的案件，因為理由在邏輯上比其所支持的結果更寬廣。法院既然已經且必然在其意見中提供建議，也就是諮詢意見，則想像它可以比至今更認真考慮如何更有效地履行這項功能，只是「the smallest of steps」〔最小的一步〕\pg{481 頁}。

\section{結論\pg{481 至 482 頁}}
\label{sec:AS-5}

\subsection{（一）Marshall 的主張與法的統一性\pg{481 至 482 頁，註 75--77}}

結論起首說，John Marshall\ix{persons}{marshall@Marshall, John（馬歇爾）} 主張司法部門的職責是「to say what the law is」〔說明法律是什麼〕（註 75 引 Marbury\ix{cases}{02@美國!01@聯邦最高法院!1803Marbury v. Madison（1803）@Marbury v. Madison（1803）} v. Madison\ix{persons}{madison@Madison, James（麥迪遜）}, 5 U.S. (1 Cranch) 137, 177 (1803)），如今流行將它視為 Marshall 式傲慢的典型，對 Marbury 一案的結果不必要，憲法未予支持，政策上也不明智\pg{481 至 482 頁，註 76 對照 Nagel\ix{persons}{nagel@Nagel, Robert F.} 稱司法至上\ix{terms}{01@\zhlidxcat{司法審查權限之歸屬}!司法至上}的主張為「self-importance」〔自重〕的表現}。

兩人的答覆是：若有多個機構說明法律是什麼，就可能有多部法律，更準確地說，對同一部法律有多種解釋\pg{482 頁}。若如 Fuller\ix{persons}{fuller@Fuller, Lon L.（富勒）} 所主張，知道法律是什麼與知道如何遵從是法治本身的必要條件（註 77 引 Fuller, \en{The Morality of Law}），則繁多的法律與繁多的解釋與法治本身相悖。他們承認意見的多樣是有價值的社會現象，但法律的存在，主要是因為意見與行動的多樣有時產生的害處大於好處；當重視協調與定紛\ix{terms}{02@\zhlidxcat{定紛止爭論}!定紛止爭}多於多樣性時，就使用法律。Marshall 的主張在兩人看來「nothing less than」〔不外乎〕後來由 Fuller 精煉的這項觀察：沒有單一而有權威的解釋者，法律與哲學書籍及建議專欄裡許多不被執行的指示之間，將差別甚微\pg{482 頁}。

\subsection{（二）讓步與結語\pg{482 頁，註 78}}

兩人承認，Marshall\ix{persons}{marshall@Marshall, John（馬歇爾）} 作為法官，從需要單一權威解釋者過快地推到司法部門是唯一能擔任此角色的機構。其他機構雖然是候選者，但沒有一個具備能使其以與最高法院同等程度的一致性發言的制度性約束。只要如此，而且沒有證據顯示並非如此，則說明法律是什麼就將、也應是最高法院的功能\pg{482 頁}。註 78 指出，與批評者交換意見最有價值的成果之一，是使辯論的核心，至少是兩人與主張由國會擔任單一權威解釋者的人之間的核心，經驗問題的面貌更為清楚；如 Peabody\ix{persons}{peabody@Peabody, Bruce G.} 所觀察，這場辯論可能成為嚴肅經驗研究的跳板，而這種研究尚未進行\pg{註 78 引 Peabody 85--90 頁}。

最後一段區分兩種位置。對向最高法院論辯、或評價其工作的人，最高法院所說與其所稱解釋的法律之間，始終存在某種概念空間。對最高法院向其發言的人，「the functions of law have been better served, and will be better served in the future, if the Constitution is what the Supreme Court says it is」〔如果憲法就是最高法院所說的那樣，法律的功能過去得到了更好的實現，將來也會得到更好的實現〕\pg{482 頁}。
